\chapter{Data Science}

Data Science is a multidisciplinary field that combines statistical analysis, machine learning, and domain expertise to extract insights and knowledge from data. 

\section{Statistics}
Data science involves, in a large scale, a lot of statistics, which is the branch of mathematics that deals with the collection, analysis, interpretation, presentation, and organization of data. Some of the most important concepts in statistics include:
\begin{itemize}
    \item \ul{Outlier}: A data point that is significantly different from the other data points in the dataset. Outliers can be caused by errors in data collection or by natural variability in the data. They can have a significant impact on the results of statistical analysis, so it is important to identify and handle them appropriately.
    \item \ul{Mean}: The average of a set of numbers. It represents the central tendency of the data.
    \begin{equation*}
        \text{Mean} = E[X] = \mu := \frac{1}{n} \sum_{i=1}^{n} x_i
    \end{equation*}
    \item \ul{Median}: The middle value of a set of numbers when they are arranged in order. It is less affected by outliers than the mean. Other terms as percentiles and quartiles are also used to describe the distribution of the data. The $p$-th percentile is the value below which $p\%$ of the data falls. The first quartile (Q1) is the 25th percentile, the second quartile (Q2) is the 50th percentile (which is also the median), and the third quartile (Q3) is the 75th percentile.
    \item \ul{Mode}: The value that appears most frequently in a dataset. It is useful for categorical data.
    \item \ul{Variance}: A measure of how much the data varies from the mean. It is calculated as the average of the squared differences from the mean.
    \begin{equation*}
        \text{Variance} = E[(X - \mu)^2] = \sigma^2 := \frac{1}{n} \sum_{i=1}^{n} (x_i - \mu)^2
    \end{equation*}
    It can be however calculated using the Besel correction, which divides by $n-1$ instead of $n$ to provide an unbiased estimator of the population variance when working with a sample. This is because when we calculate the sample mean, we are using the same data points that we are trying to estimate the variance for, which can lead to an underestimation of the true population variance. By dividing by $n-1$, we are effectively increasing the variance estimate to account for this bias.
    \begin{equation*}
        \text{Sample Variance} = \frac{1}{n-1} \sum_{i=1}^{n} (x_i - \mu)^2
    \end{equation*}
    \item \ul{Standard Deviation}: The square root of the variance. It is a measure of how much the data varies from the mean, but it is in the same units as the data.
    \begin{equation*}
        \text{Standard Deviation} = \sqrt{\text{Variance}} = \sigma
    \end{equation*}
    \item \ul{Skewness}: A measure of the asymmetry of the distribution of the data. A positive skew indicates that the data is skewed to the right, which means that the mean is greater than the median, meaning that there are more values on the right side of the distribution than on the left. A negative skew indicates that the data is skewed to the left, which means that the mean is less than the median, meaning that there are more values on the left side of the distribution than on the right. A skewness of zero indicates that the data is perfectly symmetrical.
    \item \ul{Kurtosis}: A measure of the "tailedness" of the distribution of the data. A high kurtosis indicates that the data has heavy tails, which means that there are more extreme values in the data than would be expected in a normal distribution. A low kurtosis indicates that the data has light tails, which means that there are fewer extreme values in the data than would be expected in a normal distribution. A kurtosis of zero indicates that the data has the same tailedness as a normal distribution.
    \item \ul{Covariance}: A measure of how much two random variables change together. It is calculated as:
    \begin{equation*}
        \text{Cov}(X, Y) = E[(X - E[X])(Y - E[Y])]
    \end{equation*}
    A positive covariance indicates that the two variables tend to increase or decrease together, while a negative covariance indicates that one variable tends to increase when the other decreases. A covariance of zero indicates that there is no linear relationship between the two variables.
    \item \ul{Correlation}: A measure of the strength and direction of the linear relationship between two random variables. It is calculated as:
    \begin{equation*}
        \text{Correlation}(X, Y) = r_{X,Y} = \frac{\text{Cov}(X, Y)}{\sigma_X \sigma_Y}
    \end{equation*}
    where $\sigma_X$ and $\sigma_Y$ are the standard deviations of $X$ and $Y$, respectively. The correlation coefficient ranges from -1 to 1, where a value of 1 indicates a perfect positive linear relationship, a value of -1 indicates a perfect negative linear relationship, and a value of 0 indicates no linear relationship between the two variables.

    It should be noted that correlation does not imply causation. Just because two variables are correlated does not mean that one variable causes the other to change. There may be other factors at play that are influencing both variables, or the correlation may be purely coincidental. In \myhref{https://www.tylervigen.com/spurious-correlations}{this website} you can find some examples of spurious correlations, which are correlations that are not causally related.
\end{itemize}

\section{Data}
Data is the raw material that fuels the development of AI systems, and understanding how to work with data is crucial for anyone interested in the field of AI.
% Two main people work on it:
% \begin{itemize}
%     \item \ul{Data Engineer}: Responsible for building and maintaining the infrastructure that allows for the collection, storage, and processing of data. 
%     \item \ul{Data Scientist}: Responsible for analyzing and interpreting data to extract insights and knowledge.
% \end{itemize}

When working with data, the PPDAC Data Cycle is often used as a framework to guide the process. The PPDAC Data Cycle consists of five stages:
\begin{itemize}
    \item \ul{Problem}: Defining the problem that needs to be solved, stating the objectives and questions that need to be answered.
    \item \ul{Plan}: Developing a plan for how to answer the questions and achieve the objectives. This includes deciding on which data to collect, how to collect it, and how to analyze it.
    \item \ul{Data}: Collecting the data according to the plan. This can involve gathering data from various sources, such as databases, APIs, or web scraping. It also involves preparing the data for analysis, which may include cleaning, transforming, and encoding the data.
    \item \ul{Analysis}: Analyzing the data to extract insights and knowledge. This can involve using various statistical and machine learning techniques to identify patterns and relationships in the data.
    \item \ul{Conclusion}: Drawing conclusions from the analysis and communicating the results. This can involve creating visualizations, writing reports, or presenting the findings to stakeholders.
\end{itemize}

\subsection{Data Collection}
    

Regarding data collection, there are some aspects that must be taken into account:
\begin{itemize}
    \item \ul{Data Quality}: The quality of the data is crucial for the success of any AI project.
    \item \ul{Data Integrity}: Ensuring that the data is accurate and reliable.
    \item \ul{Data Consistency}: Ensuring that the data is consistent across different sources and formats.
\end{itemize}

\subsection{Data Preprocessing}


Once the data has been collected, it must be preprocessed before it can be used for training AI models. This involves:
\begin{itemize}
    \item \ul{Data Validation}: Checking the data for errors and inconsistencies. Depending on the source and type of data, some rules can be applied to validate it.
    \item \ul{Data Sorting}: Organizing the data in a way that makes it easier to work with. This can involve sorting the data by a specific column.
    \item \ul{Data Aggregation}: Combining data from multiple sources or records into a single record. This can involve calculating summary statistics, such as the mean or median, for a group of records. This is often used to reduce the size of the data and to make it easier to analyze. However, it can also lead to loss of information, so it should be used with caution.
    \item \ul{Data Cleaning}: Removing or correcting any errors or inconsistencies in the data.
    \item \ul{Data Transformation}: Converting the data into a format that can be used for training AI models.
\end{itemize}

Data preprocessing also involves data scaling, which is the process of normalizing or standardizing the data to ensure that it is on the same scale. This is important because many AI algorithms are sensitive to the scale of the data and can perform poorly if the data is not scaled properly. Some common scalers include:
% // TODO: Memorizar Scalers
\begin{itemize}
    \item \ul{\texttt{Binarizer}}: Converts categorical data into binary data. A threshold is set, and values above the threshold are converted to 1, while values below the threshold are converted to 0.
    \begin{equation*}
        X' = \begin{cases}
            0 & \text{if } X \leq \text{threshold} \\
            1 & \text{if } X > \text{threshold}
        \end{cases}
    \end{equation*}
    \item \ul{\texttt{MaxAbsScaler}}: Scales the data to the range $[-1, 1]$ by dividing each value by the maximum absolute value in the data.
    \begin{equation*}
        X' = \frac{X}{\max(|X|)}
    \end{equation*}
    \item \ul{\texttt{MinMaxScaler}}: Scales the data to a specified range, let's say $[\min, \max]$. Firsrly the data is scaled to the range $[0, 1]$ ($X_{\text{std}}$), and then it is scaled to the specified range.
    \begin{equation*}
        X_{\text{std}} = \frac{X - \min(X)}{\max(X) - \min(X)}\qquad X' = X_{\text{std}} \cdot (\max - \min) + \min
    \end{equation*}
    \item \ul{\texttt{Normalizer}}: Scales the data to have a unit norm. This is done by dividing each value by the norm of the data. Usually, the L2 norm is used.
    \begin{equation*}
        X' = \frac{X}{\|X\|}
    \end{equation*}
    \item \ul{\texttt{PowerTransformer}}: Applies a power transformation to the data to make it more Gaussian-like. This is useful for data that is not normally distributed.
    \item \ul{\texttt{RobustScaler}}: Scales the data using statistics that are robust to outliers. This is useful for data that contains outliers.
    \begin{equation*}
        X' = \frac{X - \text{median}(X)}{\text{IQR}(X)}
    \end{equation*}
    where $\text{IQR}(X)$ is the interquartile range of $X$.
    \item \ul{\texttt{StandardScaler}}: Scales the data to have a mean of 0 and a standard deviation of 1.
    \begin{equation*}
        X' = \frac{X - \mu}{\sigma}
    \end{equation*}
    where $\mu$ is the mean of $X$ and $\sigma$ is the standard deviation of $X$.
\end{itemize}

One important aspect of data preprocessing is \emph{encoding} the data. This is the process of converting categorical data into numerical data that can be used for training AI models. There are several techniques for encoding data, such as:
% // TODO: Memorizar Encodings
\begin{itemize}\label{item:encoding}
    \item \ul{\texttt{Ordinal Encoding}}: Assigning a unique integer to each category in the data. The order of the integers is important, as it implies a relationship between the categories.
    \begin{equation*}
        \left\{\text{small}, \text{medium}, \text{large}\right\} \rightarrow \left\{0, 1, 2\right\}
    \end{equation*}

    \item \ul{\texttt{Label Encoding}}: Similar to ordinal encoding, but it does not imply any relationship between the categories. Each category is assigned a unique integer, but the order of the integers does not matter.
    \begin{equation*}
        \left\{\text{red}, \text{green}, \text{blue}\right\} \rightarrow \left\{0, 1, 2\right\}
    \end{equation*}
    
    \item \ul{\texttt{One-Hot Encoding}}: Creating a binary column for each category in the data.
    \begin{equation*}
        \left\{\text{red}, \text{green}, \text{blue}\right\} \rightarrow \begin{bmatrix}
            1 & 0 & 0 \\
            0 & 1 & 0 \\
            0 & 0 & 1
        \end{bmatrix}
    \end{equation*}

    \item \ul{\texttt{Dummy Encoding (Leave out the last)}}: Similar to one-hot encoding, but it leaves out the last category to avoid multicollinearity.
    \begin{equation*}
        \left\{\text{red}, \text{green}, \text{blue}\right\} \rightarrow \begin{bmatrix}
            1 & 0 \\
            0 & 1 \\
            0 & 0
        \end{bmatrix}
    \end{equation*}

    \item \ul{\texttt{Absolute Frequency Encoding}}: Assigning a unique integer to each category based on the absolute frequency of the category in the data. Once drawback of this technique is that, if two categories have the same frequency, they will be assigned the same integer, which can lead to confusion.
    \begin{gather*}
        \left\{\text{red}, \text{red}, \text{red}, \text{green}, \text{blue}, \text{blue}\right\}\\ \left\{\text{red}, \text{green}, \text{blue}\right\} \rightarrow \left\{3, 1, 2\right\}
    \end{gather*}

    \item \ul{\texttt{Relative Frequency Encoding}}: Assigning a unique integer to each category based on the relative frequency of the category in the data.
    \begin{gather*}
        \left\{\text{red}, \text{red}, \text{red}, \text{green}, \text{blue}, \text{blue}\right\}\\ \left\{\text{red}, \text{green}, \text{blue}\right\} \rightarrow \left\{\nicefrac{3}{6}, \nicefrac{1}{6}, \nicefrac{2}{6}\right\}
    \end{gather*}

    \item \ul{\texttt{Feature Hash Encoding}}: Hashing the categories into a fixed number of columns. This is useful for data with a large number of categories.
\end{itemize}

\subsection{Data Usage}

Once the data has been preprocessed, they can be used. There are two main destinations for the data:
\begin{itemize}
    \item \ul{Data Visualization}: The process of creating visual representations of data to help understand and communicate insights. This is used by people.
    
    The Gestalt Principles of Perception are a set of principles that describe how people perceive and organize visual information, and they can be applied to data visualization to create more effective and engaging visualizations.
    \item \ul{Data Modeling}: The process of using data to train AI models. This is used by machines.
\end{itemize}